% §18.1 The Encrypted Search Problem
%
% Monograph register: the problem (def:encrypted-search), three naive
% resolutions and why each fails, the cipher-map framing, the
% leakage-recovery trade-off (1 - PPV). Tightened. Uses \Cref{def:ppv}
% for the known-good PPV label. Cross-refs def:ppv, def:cipher-map,
% sec:fhe-sidebar, sec:maximizing-confidentiality, sec:bernoulli-error-model,
% sec:bloom-based, sec:confidentiality-measures, sec:adversarial-estimation,
% sec:toy-system. DEFINED: sec:encrypted-search-problem, def:encrypted-search.

\section{The Encrypted Search Problem}
\label{sec:encrypted-search-problem}

A client stores documents on a remote server and wants to search them for
a keyword without trusting the server to read either the documents or the
queries. The documents hold private content, the queries reveal what the
client is interested in, and handing either to the server in the clear
defeats the point. The server has abundant storage and computation and
will do work on the client's behalf; what it cannot be trusted with is
reading the data it stores.

\begin{definition}[Encrypted-search problem]
\label{def:encrypted-search}
A client holds a corpus of plaintext documents $\{d_1, \ldots, d_n\}$,
each a set of keywords from a vocabulary $W$. The client wants to
outsource the corpus to a server so that, for any keyword $w \in W$, it
can issue a query the server resolves by returning the identifiers $\{i :
w \in d_i\}$, while the server learns neither the plaintext of any
document nor the keyword of any query.
\end{definition}

The two privacy goals are intertwined: whatever operation tests a query
against a stored document is the seam through which information leaks. A
scheme leaking nothing about either side does not exist in general, since
the server must learn something to return an answer, if only the bare
fact that a query matched. The design question is how much, and of what
kind.

\subsection*{Three naive resolutions}

The first resolution skips the server: download everything and search
locally. It leaks nothing but eliminates every reason to use a server,
since a client that can download and store the full corpus could have
hosted it locally. The second preserves the server at the cost of
privacy: send queries in plaintext and let the server match them. If the
documents are plaintext the server reads everything; if they are
encrypted but the queries are not, the server still learns the workload
(which keywords matter, how often, in what order) and correlates it with
access patterns, timing, and document sizes to recover content in many
cases. Encrypting the documents but not the queries protects the corpus
only as long as the workload is uninformative, and most realistic
workloads are extremely informative. The third chases full privacy
through fully homomorphic encryption
(\Cref{sec:fhe-sidebar}), which lets the server evaluate an arbitrary
circuit on encrypted inputs but is many orders of magnitude too slow for
the corpus scales encrypted search must address. The first abandons the
server, the second privacy, the third performance. The fourth resolution,
this chapter's, abandons \emph{exactness}.

\subsection*{Encrypted search as a cipher map}

The cipher map of \Cref{def:cipher-map} operates on encrypted
representations and returns the right answer up to bounded $\fprate$ and
$\fnrate$, and encrypted search fits it directly. The plaintext function
is the membership test ``does document $d$ contain keyword $w$?''; the
encrypted-domain map is a server-side test on whatever encrypted form of
$d$ and $w$ the construction provides; the error rates are the share of
non-members reported as matches and the share of true members missed. A
search is a sequence of such tests, one per document, and the Kronecker
composition law of \Cref{sec:bernoulli-error-model} aggregates the
per-test rates into a corpus-wide result. Encrypted search is thus a
family of cipher maps parameterized by the underlying approximate-set
structure and the encryption scheme applied to it; the two rates govern
recovery (how faithfully the response reflects the true match set) and
the confidentiality of \Cref{sec:maximizing-confidentiality} governs
leakage (how much the response reveals about the plaintext). The Bernoulli
machinery applies unchanged; only the operational reading of $\fprate,
\fnrate$ changes.

\subsection*{The leakage-recovery trade-off}

A tighter system (smaller $\fprate, \fnrate$) returns a more faithful
match set but also leaks more: its membership reports are more reliable,
so a third party who learns a report has a higher posterior on whether
the document truly contained the queried keyword. The positive predictive
value of \Cref{def:ppv}, $\Pr[w \in d \mid \text{match reported}]$, read
with the adversary in mind, is the leakage measure, and confidentiality
is $1 - \PPV$ (\Cref{sec:maximizing-confidentiality}), the two halves of
the trade-off being opposite faces of one formula. The trade-off is
unavoidable: an exact system has $\PPV = 1$ and confidentiality zero, a
system returning random matches has confidentiality $1 - \pi$ and is
useless, and every construction picks a point between.
\Cref{sec:bloom-based} builds two Bloom-filter constructions,
\Cref{sec:confidentiality-measures} works out the PPV/NPV/F1 measures as
confidentiality bounds, \Cref{sec:adversarial-estimation} addresses
estimating leakage under adversarial workloads, and
\Cref{sec:toy-system} walks an end-to-end implementation.
